\chapter{Implementierung}
\label{sec:implementierung}
In diesem Kapitel wird die Umsetzung der geplanten Softwarearchitektur im Detail beschrieben. Der Fokus liegt dabei auf der Backend-Implementierung, während das Frontend nur oberflächlich betrachtet wird, dessen Hauptfunktion im Anzeigen und Bedienen der Applikation besteht. 

Im folgenden wird zunächst die Gesamtarchitektur betrachtet, wobei Besonderheiten der Implementierung sowie die Nutzung von Django hervorgehoben werden. Anschließend wird die eigens entwickelte Django App \textit{Voice} vorgestellt, in der die Geschäftslogik implementiert ist. Zuletzt wird auf den Aufbau des eigentlichen LLM-Agenten eingegangen.

\section{Proof on Concept}
Zeigen des ersten Prototypen welcher Mit jupyter Notebook entstanden ist
Probleme die sich ergeben haben (Threading) 

\section{Backendstruktur}
Ganz kurz beschreiben wie Struktur wie Notisent backend aufgebaut ist. Daran anschließend erklären wie eigenes Projekt eingebunden wird

\section{Django voice-App}
Verschiedene Apps, Django 
DsPy 
Webhooks 
Code wie der Prompt eingebunden ist

\section{Teams Bot}
Attendee API 
Wie funktioniert, wie wurde eingesetzt
Vorherige Probleme mit erstellen von Teams Bot über MS Graph API und Azure erwähnen

\section{Verarbeitung des Live-Transkripts}
Prompt 

\section{Frontend Application}
Weniger Fokus drauf, nur anzeige von Zusammengefassten Texten, TBC 

\section{Optimierung}
Doppeloptimierung, Sprachmodell zum erkennen 
Problem der Zugänglichkeit zu MS Team liveaudio
Anpassen der Darstellung von 
zero shot prompting 
frew shot prompting 
wie kann das ergebnis des assistenten optimiert werden
evtl voranalyse? SpaCy
